% 西南科技大学研究生学位论文主文件
% 编译顺序：XeLaTeX → BibTeX → XeLaTeX → XeLaTeX。
\documentclass[UTF8,a4paper,twoside,openany,zihao=-4,fontset=none]{ctexbook}

% 只改这一行选择模板：professional-master / academic-master / academic-doctor
\providecommand{\ThesisVariant}{professional-master}
\usepackage{swust-thesis}

% ===== 填写论文信息 =====
\renewcommand{\ThesisTitle}{论文题目}
\renewcommand{\ThesisTitleEN}{\EnglishTitlePlaceholder} % 填写正式英文题目时替换右侧内容
\renewcommand{\AuthorName}{XXX}
\renewcommand{\AuthorNameEN}{XXX}
\renewcommand{\StudentID}{XXXXXXXXXXX}
\renewcommand{\Supervisor}{XXX 教授}
\renewcommand{\SupervisorEN}{XXX}
\renewcommand{\ExternalSupervisor}{} % 没有校外导师则留空
\renewcommand{\College}{XXXX学院}
\renewcommand{\DefenseDate}{20XX年XX月XX日}
\renewcommand{\ThesisDate}{20XX年XX月}
\renewcommand{\ThesisDateEN}{Month 20XX}
\renewcommand{\CommitteeMembers}{XXX、XXX、XXX}
\renewcommand{\Reviewers}{XXX、XXX}
\renewcommand{\Classification}{}
\renewcommand{\UDCCode}{}
\renewcommand{\SecurityLevel}{}
\renewcommand{\KeywordsCN}{关键词一；关键词二；关键词三}
\renewcommand{\KeywordsEN}{Keyword one; Keyword two; Keyword three}
\renewcommand{\ThesisTypeCN}{应用研究}
\renewcommand{\ThesisTypeEN}{Application Research}

% 学位类型和学科/领域随模板选择有示例默认值；按实际情况修改时取消注释。
% \renewcommand{\DegreeType}{实际学位类型}
% \renewcommand{\DegreeTypeEN}{Official English Degree Name}
% \renewcommand{\MajorField}{实际学科或领域}

% 长题目与名单按需要分成两行；不要在 ThesisTitle 中手动插入换行。
% \renewcommand{\CoverTitleFirstLine}{题目第一行}
% \renewcommand{\CoverTitleSecondLine}{题目第二行}
% \renewcommand{\CommitteeMembersSecondLine}{其余答辩委员}

\begin{document}

\SWUSTMakeCover
\SWUSTMakeTitlePages
\SWUSTMakeDeclarations

\SWUSTStartChineseAbstract
% 中文摘要：硕士约 600--800 字并突出新见解；博士约 1000--1500 字并突出创新点。
请在此填写中文摘要，约 \AbstractLengthCN{} 个汉字，突出论文的\AbstractFocusCN。摘要正文每段首行缩进两个汉字，不放图、表或非公知符号；段间不另空行。

\SWUSTEndChineseAbstract

\SWUSTStartEnglishAbstract
Please enter the English abstract here. It should accurately correspond to the Chinese abstract. English paragraphs are flush left and are not separated by blank lines.

\SWUSTEndEnglishAbstract

\SWUSTMakeTOC
% 若需要主要符号表，可在这里添加：\SWUSTUnnumberedChapter{主要符号表}

\SWUSTStartMain
\chapter{绪论}
\section{研究背景与意义}
请在此填写研究背景、研究价值与待解决的问题。正文使用小四号字、1.2 倍行距，段前段后不额外留白。本模板的版式依据学校 2024 年版撰写规范\upcite{swust2024}。此处演示顺序编码引用，写作时请替换为研究中实际使用的资料。

\subsection{研究现状}
请在此归纳国内外研究现状，并说明本文的研究思路。注释示例\footnote{注释用于解释性、说明性或补充性材料，并按页重新编号。}显示在当前页底部。

\chapter{研究内容与方法}
\section{研究方案}
图应先在正文提及，再紧随相关文字放置。图~\ref{fig:example} 是图题与分章编号示例。

\begin{figure}[htbp]
  \centering
  \fbox{\parbox[c][5cm][c]{6.67cm}{\centering 请在此放入清晰的图片}}
  \caption{示例图题}\label{fig:example}
\end{figure}

表~\ref{tab:example} 为三线表示例。表题在表格上方，表与正文同宽，数字按列对齐。
\begin{table}[htbp]
  \centering
  \caption{示例数据表}\label{tab:example}
  \zihao{5}
  \begin{tabular*}{\textwidth}{@{\extracolsep{\fill}}lcc@{}}
    \toprule[1.5pt]
    项目 & 数值 / 单位 & 说明 \\
    \midrule[1pt]
    样本 A & 10 & 示例数据 \\
    样本 B & 12 & 示例数据 \\
    \bottomrule[1.5pt]
  \end{tabular*}
\end{table}

公式单独成行，按章编号为括号中的“章号-序号”。
\begin{equation}\label{eq:example}
  y=ax+b
\end{equation}
\noindent 式中，$y$ 为响应量；$x$ 为自变量；$a$、$b$ 为参数。

\chapter{结论}
在此概括主要结论、新见解、实际价值与进一步研究问题。


\SWUSTUnnumberedChapter{致\quad 谢}
在此感谢导师及对论文有直接贡献的个人和单位。致谢不超过 1000 个汉字。


% 仅列出正文引用的文献，按首次引用顺序编号。
\bibliographystyle{gbt7714-numerical}
\bibliography{references}

% 无附录时可注释下两行。
\SWUSTStartAppendices
\chapter{补充材料}
在此放置正文中不宜展开的推导、问卷、辅助数据或程序。附录的图、表和公式与正文分别编号。


\SWUSTUnnumberedChapter{攻读学位期间取得的研究成果}
按参考文献格式，连续编号列出与本学位论文相关的已发表/录用论文、专著、授权专利、科研获奖等；没有相关成果时按学院要求处理。


\end{document}
